\documentclass[../../../include/open-logic-section]{subfiles}

\begin{document}

\olfileid{sth}{z}{milestone}
\olsection{$\Z^-$: يو مهم پړاو}

په راتلونکې برخه کښې به \stagesinf{} ته بيا ستانۀ شو.
خو د نامتناهيت بديهي اصل په لرلو سره يو مهم پړاو ته رسېدلي يو.
اوس د تيورۍ $\Zminus$ دپاره ټول اړين بديهي اصلونه لرو.
په تفصيل سره:

\begin{defn}
تيوري $\Zminus$ دا بديهي اصلونه لري: د غړو له مخې برابري،
اتحاد، جوړې، د فرعي سټونو سټونه، نامتناهيت، او د بېلونې د
شېما ټولې بېلګې.
\end{defn}

دا نوم د \emph{زرمېلو} د سټ تيورۍ ښکارندويي کوي (يو شے ترې
\emph{کم} دے چې وروسته به ورته راشو). دا وياړ د زرمېلو حق دے،
ځکه چې همدا تيوري ئې په اساسي توګه په خپل
\citeyear{Zermelo1908Untersuchungen} کښې جوړه کړې وه.\footnote{د تاريخ
او تخنيکي جزئياتو د په زړه پورې بحث دپاره \citet[Appendix
A]{Potter2004} وګورئ.}

دا تيوري دومره پياوړې ده چې ډېر زيات رياضيکي کار پکښې کولے شو.
په ځانګړې توګه، \emph{پکار ده} چې \olref[sfr][][]{part} ته بېرته
وګورئ او ځان قانع کړئ چې ټول هغه څه چې په ساده توګه مو کړي وو،
په لا رسمي ډول په~$\Zminus$ کښې کېدلے شي. (يو څه د دې کار تر
کولو وروسته ښايي وغواړئ چې مخته لاړ شئ او
\olref[sth][z][arbintersections]{sec} ولولئ.) نو له دې وروسته،
بې له نورو څرګندونو، ځانونه په $\Zminus$ کښې (لږ تر لږه) کار کوونکي ګڼو.

\end{document}
